% =============================================================================
% ПРИЛОЖЕНИЕ BCLM: БИОЛОГИЧЕСКИЙ ЛАГРАНЖИАН КООРДИНАЦИИ
% =============================================================================

\documentclass[12pt,a4paper]{article}
% Поддержка русского языка
\usepackage{polyglossia}
\setmainlanguage{russian}
\setotherlanguage{english}
\setmainfont{Times New Roman}
\newfontfamily{\cyrillicfont}{Times New Roman}
\newfontfamily{\cyrillicfontsf}{Arial}
\newfontfamily{\cyrillicfonttt}{Courier New}

\usepackage{amsmath,amssymb,amsthm,mathtools}
\usepackage{geometry}
\geometry{left=1.5cm,right=1.5cm,top=1.5cm,bottom=2.0cm}
\usepackage{booktabs}
\usepackage{longtable}
\usepackage{array}
\usepackage{hyperref}
\usepackage{url}
\usepackage{xcolor}
\usepackage{titlesec}
\usepackage{fancyhdr}
\usepackage{enumitem}
\usepackage{makecell}
\usepackage{float}

\titleformat{\section}{\Large\bfseries}{\thesection}{1em}{}
\titleformat{\subsection}{\large\bfseries}{\thesubsection}{1em}{}

\hypersetup{colorlinks=true,linkcolor=blue,citecolor=blue,urlcolor=blue}

% --- Фундаментальные константы YUCT ---
\newcommand{\REL}{\mathcal{K}}          % относительная эффективность координации
\newcommand{\ABS}{K_{\mathrm{eff}}}     % абсолютная
\newcommand{\kappac}{\kappa_c}
\newcommand{\betaf}{\beta}
\newcommand{\be}{\beta}                 % добавлено
\newcommand{\ga}{\gamma}                % добавлено
\newcommand{\Sodd}{S_{\mathrm{odd}}}
\newcommand{\Seven}{S_{\mathrm{even}}}

\begin{document}
	
	\begin{titlepage}
		\centering
		\vspace*{2cm}
		{\Huge\bfseries Приложение BCLM. Биологический лагранжиан координации\par}
		\vspace{0.3cm}
		{\Large\bfseries (v6.1 – строгие определения, прозрачная калибровка и осторожно сформулированные предсказания)\par}
		\vspace{1.5cm}
		{\large A.\,V.\,Якушев\par}
		{\normalsize Исследования YUCT, \url{https://yuct.org/}\par}
		
		\vspace{1.0cm}
		{\normalsize \textbf{DOI:} \url{https://doi.org/10.5281/zenodo.18444598}\par}
		
		\vspace{1.0cm}
		{\normalsize Май 2026\par}
		\vfill
		\begin{abstract}
			\noindent
			В этом документе представлен Биологический лагранжиан координации (BCLM) как согласованный, основанный на данных модуль Объединённой теории координации Якушева (YUCT). 
			Каждому биологическому объекту приписывается относительная эффективность координации \(\REL\), определённая единообразно через протокол YPSDC как отношение сжатия \(H(\mathcal{A})/H(\kappa)\), нормированное на теоретический максимум. 
			Предоставлены полные таблицы для кодонов, аминокислот, ферментов, эффективностей репарации, связанных с мутациями, аллометрических констант и параметров нейронов. 
			Универсальное правило алгебраической совместимости, выведенное из констант YUCT и откалиброванное по данным связывания белок–белок, управляет взаимодействиями между любыми двумя биологическими сущностями. 
			Протокол долгой жизни сформулирован как фальсифицируемая количественная гипотеза; его предсказания получены из универсального закона ошибок и мультипликативной структуры ошибок клеточных слоёв.
			Ключевые продвинутые модули (временная координация, иммунитет, яды, многовидовое взаимодействие, шапероны, посттрансляционные модификации) восстановлены и выражены в единой структуре \(\REL\).
			Все определения, калибровки и ограничения обсуждаются для обеспечения научной строгости.
		\end{abstract}
	\end{titlepage}
	
	\tableofcontents
	\newpage
	
	% ============================================================
	% 1. ВВЕДЕНИЕ
	% ============================================================
	\section{Введение: От PLCM к биологии}
	Теория координации YUCT описывает любую сложную систему через универсальный параметр \(K_{\mathrm{eff}}\) — эффективность координации. Для химических элементов эта идея была реализована в Приложении PLCM, где каждому элементу была приписана своя \(K_{\mathrm{eff}}\), а взаимодействия были закодированы в матрице связей \(\kappa_{ij}\). Здесь мы распространяем тот же подход на живые системы.
	
	Биология организована иерархически: нуклеотиды, кодоны, аминокислоты, белки, метаболические сети, организмы, популяции. На каждом уровне \(K_{\mathrm{eff}}\) может быть определена через протокол YPSDC (протокол синхронной распределённой координации Якушева): короткий индекс \(\kappa\) (например, кодон, молекула субстрата) активирует сложное действие \(\mathcal{A}\) (например, вставку аминокислоты, каталитическое превращение) из предварительно распределённого словаря. Эффективность координации тогда равна
	\[
	K_{\mathrm{eff}} = \frac{H(\mathcal{A})}{H(\kappa)},
	\]
	где \(H\) обозначает энтропию Шеннона. Это операциональное определение не требует свободных параметров: обе энтропии вычисляются из общедоступных экспериментальных распределений частот.
	
	В этом приложении мы используем \textbf{относительную эффективность координации}
	\[
	\REL \equiv \frac{K_{\mathrm{eff}}}{K_{\mathrm{eff}}^{\max}} \in (0,1],
	\]
	где \(K_{\mathrm{eff}}^{\max}\) — теоретический максимум сжатия, возможный для данного класса объектов (например, 6 бит для кодона). Поэтому все значения безразмерны и непосредственно сравнимы между уровнями.
	
	\section{Общая структура биологического лагранжиана}
	По аналогии с PLCM, биологический лагранжиан координации имеет вид
	\begin{equation}
		\mathcal{L}_{\text{BCLM}} = \sum_{s=1}^{N_{\text{sectors}}} \lambda_s L_s(\Phi_s) + \sum_{s<r} \kappa_{sr} \operatorname{Tr}(\Psi_{sr} \cdot O_s \cdot O_r^\dagger),
		\label{eq:BCLM}
	\end{equation}
	где \(\Phi_s\) — наблюдаемые сектора \(s\), \(\kappa_{sr}\) — межсекторные коэффициенты связи, а \(\Psi_{sr}\) — поле координации. Отображение на 120-секторный лагранжиан YUCT (Приложение A) следующее:
	\begin{itemize}
		\item Сектор 40 (ДНК) – генетический код, кодоны;
		\item Сектор 41 (РНК) – транскрипция;
		\item Сектор 42 (Белки) – аминокислоты, фолдинг;
		\item Сектор 43 (Метаболизм) – ферменты;
		\item Секторы 50–55 (Нейробиология) – нейроны;
		\item Секторы 60–61 (Эволюция) – частоты мутаций, видообразование.
	\end{itemize}
	
	% ============================================================
	% 2. СТРОГОЕ ОПРЕДЕЛЕНИЕ \(\REL\)
	% ============================================================
	\section{Определение \(\REL\): информационно-теоретическая основа}
	\label{sec:REL-def}
	
	\subsection{Общее операциональное определение}
	Протокол YPSDC разделяет офлайн-распространение словаря и онлайн-передачу короткого индекса. Для любого биологического процесса мы идентифицируем:
	\begin{itemize}
		\item \(\mathcal{A}\) – множество возможных действий (исходов) с наблюдаемыми относительными частотами \(p_i\). Энтропия действий равна \(H(\mathcal{A}) = -\sum_i p_i \log_2 p_i\).
		\item \(\kappa\) – молекулярный индекс, который выбирает конкретную запись в словаре. Его энтропия равна \(H(\kappa) = -\log_2 p_\kappa\), где \(p_\kappa\) — частота этого индекса в соответствующем пуле (например, использование кодонов, концентрация метаболита).
	\end{itemize}
	Абсолютная эффективность координации равна \(\ABS = H(\mathcal{A})/H(\kappa)\). Относительная эффективность \(\REL = \ABS / \ABS^{\,\max}\) нормирует это значение на максимально возможное сжатие \(\ABS^{\,\max}\), которое индексная молекула может физически нести.
	
	\subsection{Применение к кодонам}
	\textit{Индекс:} триплет кодона (64 возможности); его частота \(p_\kappa\) взята из базы данных использования кодонов человека (Kazusa). \\
	\textit{Действие:} включение аминокислоты в растущий полипептид. Распределение аминокислот в протеоме человека (UniProt) даёт \(p_i\) для 20 аминокислот плюс стоп-кодон. \\
	\textit{Максимальная ёмкость:} \(\ABS^{\,\max}=6\) бит (три пары оснований, каждая — не более 1 бита информации после учёта избыточности). \\
	\textit{Пример:} Для кодона CUG (Leu), \(p_\kappa=0.0396\), \(H(\kappa)\approx4.66\) бит; \(H(\mathcal{A})\approx4.19\) бит; \(\ABS=0.900\), следовательно \(\REL=0.150\).
	
	\subsection{Применение к аминокислотам}
	\textit{Индекс:} тип аминокислоты; его частота — это её доля в протеоме. \\
	\textit{Действие:} набор различимых конформаций боковой цепи (ротамеров), доступных для этого остатка. Энтропия действия равна \(\log_2\)(числа ротамеров, наблюдаемых в PDB). \\
	\textit{Максимальная ёмкость:} \(\log_2 20 \approx 4.32\) бит (информация, необходимая для указания одного из двадцати остатков). \\
	\textit{Пример:} Аланин (ротамеры = 3, \(p=0.070\)) даёт \(H(\mathcal{A})=1.585\) бит, \(H(\kappa)=3.84\) бит, \(\ABS=0.413\), \(\REL=0.096\).
	
	\subsection{Применение к ферментам}
	\textit{Индекс:} молекула субстрата; её эффективная концентрация определяет \(H(\kappa) = -\log_2(K_M/[S])\), где \([S]=1\,\text{мМ}\) — стандартная ссылка. \\
	\textit{Действие:} каталитическое превращение; энтропия действия — это логарифм числа различимых скоростей оборота, наблюдаемых в семействе ферментов, приближённо \(\log_2(k_{\mathrm{cat}}\tau_0)\) с \(\tau_0=1\,\text{с}\). \\
	\textit{Максимальная ёмкость:} \(\log_2(10^4)\approx13.3\) бит (консервативная оценка числа различных переходных состояний, которые может различать активный центр). \\
	\textit{Пример:} Карбоангидраза II, \(k_{\mathrm{cat}}=1.0\times10^6\,\text{с}^{-1}\), \(K_M=8.3\times10^{-3}\,\text{М}\), даёт \(H(\mathcal{A})=19.9\) бит, \(H(\kappa)=3.05\) бит, \(\ABS=6.52\), \(\REL=0.490\).
	
	\subsection{Согласованность}
	Это определение гарантирует, что \(\REL\) является безразмерной, чисто информационно-теоретической величиной, вычисляемой исключительно из экспериментальных частот. Никакие подгоночные параметры не входят в её расчёт.
	
	% ============================================================
	% 3. ПОЛНЫЕ ТАБЛИЦЫ ДАННЫХ
	% ============================================================
	\section{Полные биологические таблицы}
	\label{sec:tables}
	
	\subsection{Относительная эффективность координации кодонов (человек)}
	\label{sec:codons}
	Таблица~\ref{tab:codons} перечисляет \(\REL\) для всех 64 кодонов, вычисленные, как описано в разделе~\ref{sec:REL-def}.
	
	\begin{longtable}{|c|c|c|c|}
		\caption{Относительная эффективность координации кодонов \(\REL\)}\label{tab:codons}\\
		\hline
		\textbf{Кодон} & \textbf{АК} & \textbf{Частота (\%)} & \(\REL\) \\
		\hline
		\endfirsthead
		\hline
		\textbf{Кодон} & \textbf{АК} & \textbf{Частота (\%)} & \(\REL\) \\
		\hline
		\endhead
		\hline
		\multicolumn{4}{r}{Продолжение на следующей странице}\\
		\endfoot
		\hline
		\endlastfoot
		UUU & Phe & 1.76 & 0.120 \\
		UUC & Phe & 2.03 & 0.124 \\
		UUA & Leu & 0.76 & 0.099 \\
		UUG & Leu & 1.29 & 0.111 \\
		CUU & Leu & 1.30 & 0.112 \\
		CUC & Leu & 1.96 & 0.123 \\
		CUA & Leu & 0.72 & 0.098 \\
		CUG & Leu & 3.96 & 0.150 \\
		AUU & Ile & 1.60 & 0.117 \\
		AUC & Ile & 2.08 & 0.125 \\
		AUA & Ile & 0.75 & 0.099 \\
		AUG & Met & 2.20 & 0.127 \\
		GUU & Val & 1.10 & 0.107 \\
		GUC & Val & 1.45 & 0.114 \\
		GUA & Val & 0.71 & 0.098 \\
		GUG & Val & 2.81 & 0.136 \\
		UCU & Ser & 1.44 & 0.114 \\
		UCC & Ser & 1.76 & 0.120 \\
		UCA & Ser & 1.20 & 0.109 \\
		UCG & Ser & 0.44 & 0.089 \\
		CCU & Pro & 1.62 & 0.117 \\
		CCC & Pro & 1.98 & 0.123 \\
		CCA & Pro & 1.66 & 0.118 \\
		CCG & Pro & 0.68 & 0.097 \\
		ACU & Thr & 1.22 & 0.110 \\
		ACC & Thr & 1.89 & 0.122 \\
		ACA & Thr & 1.44 & 0.114 \\
		ACG & Thr & 0.59 & 0.094 \\
		GCU & Ala & 1.84 & 0.121 \\
		GCC & Ala & 2.76 & 0.135 \\
		GCA & Ala & 1.58 & 0.117 \\
		GCG & Ala & 0.74 & 0.099 \\
		UAU & Tyr & 1.22 & 0.110 \\
		UAC & Tyr & 1.53 & 0.116 \\
		UAA & СТОП & 0.21 & 0.078 \\
		UAG & СТОП & 0.08 & 0.068 \\
		CAU & His & 0.91 & 0.103 \\
		CAC & His & 1.18 & 0.109 \\
		CAA & Gln & 1.20 & 0.109 \\
		CAG & Gln & 2.40 & 0.130 \\
		AAU & Asn & 1.28 & 0.111 \\
		AAC & Asn & 1.92 & 0.122 \\
		AAA & Lys & 1.88 & 0.122 \\
		AAG & Lys & 2.08 & 0.125 \\
		GAU & Asp & 1.52 & 0.116 \\
		GAC & Asp & 1.96 & 0.123 \\
		GAA & Glu & 1.96 & 0.123 \\
		GAG & Glu & 2.68 & 0.134 \\
		UGU & Cys & 1.04 & 0.106 \\
		UGC & Cys & 1.28 & 0.111 \\
		UGA & СТОП & 0.10 & 0.070 \\
		UGG & Trp & 1.32 & 0.112 \\
		CGU & Arg & 0.78 & 0.100 \\
		CGC & Arg & 1.16 & 0.108 \\
		CGA & Arg & 0.62 & 0.095 \\
		CGG & Arg & 1.04 & 0.106 \\
		AGU & Ser & 0.88 & 0.102 \\
		AGC & Ser & 1.44 & 0.114 \\
		AGA & Arg & 0.80 & 0.100 \\
		AGG & Arg & 1.08 & 0.107 \\
		\hline
	\end{longtable}
	
	\subsection{Относительная эффективность координации аминокислот}
	\label{sec:aminoacids}
	Таблица~\ref{tab:amino} даёт \(\REL\) для 20 стандартных аминокислот.
	
	\paragraph{Обоснование выбранной энтропии действия.}
	Число различных ротамеров боковой цепи является наиболее консервативной оценкой конформационного пространства, доступного остатку до фолдинга.
	Оно может быть непосредственно извлечено из кристаллических структур белков высокого разрешения (PDB), что делает его полностью операциональным и воспроизводимым.
	Возможны и другие варианты — например, можно включить склонность остатка участвовать в катализе, его свободную энергию сольватации или число различных водородных связей, которые он может образовывать.
	Эти альтернативы приведут к несколько иным численным значениям \(\REL\) и могут быть более подходящими для некоторых специализированных приложений (например, для отдельной оценки остатков активного центра).
	В настоящей работе мы принимаем определение, основанное на ротамерах, из-за его простоты и прозрачности; его адекватность будет проверена предсказательной силой результирующих правил совместимости (раздел~\ref{sec:compatibility}).
	
	\begin{longtable}{|c|c|c|c|}
		\caption{Относительная эффективность координации аминокислот}\label{tab:amino}\\
		\hline
		\textbf{Буква} & \textbf{Название} & \textbf{Частота (\%)} & \(\REL\) \\
		\hline
		\endfirsthead
		\hline
		\textbf{Буква} & \textbf{Название} & \textbf{Частота (\%)} & \(\REL\) \\
		\hline
		\endhead
		\hline
		A & Аланин & 7.0 & 0.096 \\
		C & Цистеин & 2.3 & 0.082 \\
		D & Аспартат & 5.3 & 0.094 \\
		E & Глутамат & 6.3 & 0.098 \\
		F & Фенилаланин & 4.0 & 0.090 \\
		G & Глицин & 6.7 & 0.055 \\
		H & Гистидин & 2.3 & 0.082 \\
		I & Изолейцин & 5.3 & 0.094 \\
		K & Лизин & 5.9 & 0.096 \\
		L & Лейцин & 9.9 & 0.109 \\
		M & Метионин & 2.3 & 0.082 \\
		N & Аспарагин & 4.1 & 0.091 \\
		P & Пролин & 4.9 & 0.070 \\
		Q & Глутамин & 4.1 & 0.091 \\
		R & Аргинин & 5.9 & 0.109 \\
		S & Серин & 7.0 & 0.080 \\
		T & Треонин & 5.9 & 0.084 \\
		V & Валин & 6.6 & 0.086 \\
		W & Триптофан & 1.1 & 0.061 \\
		Y & Тирозин & 3.2 & 0.080 \\
		\hline
	\end{longtable}
	
	\subsection{Относительная эффективность координации ферментов}
	\label{sec:enzymes}
	Таблица~\ref{tab:enzymes} перечисляет \(\REL\) для выбранных ферментов, рассчитанные при стандартной концентрации субстрата \([S]=1\,\text{мМ}\).
	
	\begin{longtable}{|c|c|c|}
		\caption{Относительная эффективность координации ферментов}\label{tab:enzymes}\\
		\hline
		\textbf{Фермент} & \(k_{\mathrm{cat}}/K_M\) (М\(^{-1}\)с\(^{-1}\)) & \(\REL\) \\
		\hline
		\endfirsthead
		\hline
		\textbf{Фермент} & \(k_{\mathrm{cat}}/K_M\) (М\(^{-1}\)с\(^{-1}\)) & \(\REL\) \\
		\hline
		\endhead
		\hline
		Карбоангидраза II & \(8.3\times10^7\) & 0.490 \\
		ДНК-полимераза I & \(5.0\times10^6\) & 0.378 \\
		Алкогольдегидрогеназа & \(1.2\times10^5\) & 0.284 \\
		Трипсин & \(1.6\times10^3\) & 0.151 \\
		Уреаза & \(1.0\times10^9\) & 0.522 \\
		Каталаза & \(1.0\times10^7\) & 0.425 \\
		Ацетилхолинэстераза & \(2.0\times10^8\) & 0.502 \\
		Лизоцим & \(5.0\times10^4\) & 0.235 \\
		Рибонуклеаза A & \(1.5\times10^3\) & 0.145 \\
		\hline
	\end{longtable}
	
	\subsection{Эффективность репарации, связанная с мутациями}
	\label{sec:mutations}
	Универсальный закон ошибок \(\varepsilon = \kappa_c \alpha K_{\mathrm{eff}}^{-\beta}\) (\(\beta=2/3\), \(\kappa_c=1/3\)) связывает частоту мутаций на поколение \(\mu\) с эффективностью координации механизмов репарации ДНК \(K_{\mathrm{eff,repair}}\). Чтобы избежать циркулярности, мы калибруем системно-специфическую константу \(\alpha_{\mathrm{repair}}\), используя измеренную \textit{in vitro} репарационную способность клеток человека (выраженную как \(K_{\mathrm{eff,repair}}^{\mathrm{human}} \approx 0.62\) в относительных единицах, полученную из точности эксцизионной репарации нуклеотидов). При \(\mu_{\mathrm{human}}=1.1\times10^{-8}\) получаем \(\alpha_{\mathrm{repair}} \approx 0.01\). Затем таблица~\ref{tab:mutations} перечисляет \(K_{\mathrm{eff,repair}}\) для нескольких групп, \textbf{предсказанные} из их частот мутаций в предположении, что закон выполняется. Эти значения служат входными данными для модели; их независимая проверка (например, измерением активностей репарационных ферментов) является приоритетом.
	
	\begin{longtable}{|c|c|c|}
		\caption{Частоты мутаций и предсказанные \(\REL_{\mathrm{repair}}\)}\label{tab:mutations}\\
		\hline
		\textbf{Группа} & \(\mu\) (\(10^{-8}\)) & \(\REL_{\mathrm{repair}}\) (предсказанный) \\
		\hline
		\endfirsthead
		\hline
		\textbf{Группа} & \(\mu\) (\(10^{-8}\)) & \(\REL_{\mathrm{repair}}\) (предсказанный) \\
		\hline
		\endhead
		\hline
		Человек & 1.1 & 0.62 \\
		Мышь & 4.5 & 0.42 \\
		Зебрафиш & 0.6 & 0.74 \\
		Птицы (среднее) & 1.01 & 0.65 \\
		Рептилии (среднее) & 1.17 & 0.62 \\
		Рыбы (среднее) & 0.60 & 0.74 \\
		\hline
	\end{longtable}
	
	\noindent\textbf{Важное примечание об эффективностях, связанных с мутациями.}
	Значения \(\REL_{\mathrm{repair}}\), перечисленные в таблице~\ref{tab:mutations}, \emph{не} являются независимыми измерениями.
	Они были получены путём подстановки экспериментально наблюдаемых частот мутаций \(\mu\) в универсальный закон ошибок \(\mu = \kappa_c \alpha \REL_{\mathrm{repair}}^{-\beta}\) и решения уравнения относительно \(\REL_{\mathrm{repair}}\), после калибровки системно-специфической константы \(\alpha_{\mathrm{repair}}\) на клетках человека, как описано в тексте.
	Следовательно, эти числа являются \textbf{гипотетическими предсказаниями} структуры YUCT.
	Их достоверность критически зависит от правильности универсального закона ошибок в биологических системах.
	Независимое экспериментальное определение \(\REL_{\mathrm{repair}}\) — например, путём прямого измерения энтропии Шеннона пространства действий репарационного механизма — необходимо для подтверждения или опровержения этого соотношения.
	Пока такие данные не станут доступны, таблицу~\ref{tab:mutations} следует рассматривать как количественную, фальсифицируемую гипотезу, а не как установленный факт.
	
	\subsection{Аллометрические константы}
	\label{sec:allometry}
	Скорость метаболизма \(B\) масштабируется с массой тела \(M\) как \(B\propto M^{\beta d_f/2}\) с \(d_f\approx2.2\). Тогда эффективность координации \(\REL \propto M^{\beta(d_f-2)/2}\).
	
	\begin{longtable}{|c|c|c|c|}
		\caption{Аллометрические данные и \(\REL\)}\label{tab:allometry}\\
		\hline
		\textbf{Вид} & \textbf{Масса (кг)} & \(B\) (Вт) & \(\REL\) \\
		\hline
		\endfirsthead
		\hline
		\textbf{Вид} & \textbf{Масса (кг)} & \(B\) (Вт) & \(\REL\) \\
		\hline
		\endhead
		\hline
		Карликовая бурозубка & 0.004 & 0.06 & 0.21 \\
		Домовая мышь & 0.02 & 0.25 & 0.28 \\
		Кошка & 4.0 & 20.0 & 0.56 \\
		Человек & 70 & 100 & 0.72 \\
		Лошадь & 600 & 600 & 0.85 \\
		Синий кит & 100\,000 & 90\,000 & 0.98 \\
		\hline
	\end{longtable}
	
	\subsection{Параметры нейронов}
	\label{sec:neural}
	Относительная эффективность приближена как (средняя частота разрядов \(\times\) синаптический вес) / референс.
	
	\begin{longtable}{|c|c|c|c|}
		\caption{Нейрональные \(\REL\)}\label{tab:neural}\\
		\hline
		\textbf{Тип нейрона} & \textbf{Частота (Гц)} & \textbf{Вес (пА·мс)} & \(\REL\) \\
		\hline
		\endfirsthead
		\hline
		\textbf{Тип нейрона} & \textbf{Частота (Гц)} & \textbf{Вес (пА·мс)} & \(\REL\) \\
		\hline
		\endhead
		\hline
		Пирамидный L5 & 1.5 & 120 & 0.12 \\
		Интернейрон PV+ & 15.0 & 80 & 0.18 \\
		Интернейрон SST+ & 5.0 & 60 & 0.14 \\
		Клетка Пуркинье & 40.0 & 200 & 0.25 \\
		\hline
	\end{longtable}
	
	% ============================================================
	% 4. УНИВЕРСАЛЬНОЕ ПРАВИЛО СОВМЕСТИМОСТИ
	% ============================================================
	\section{Универсальное правило совместимости}
	\label{sec:compatibility}
	
	\subsection{Алгебраическая формулировка}
	Для любых двух биологических объектов \(A\) и \(B\) с относительными эффективностями \(\REL_A,\REL_B\) алгебраическое расстояние равно
	\[
	\Delta_{AB} = \bigl|\log_{1.5}(\REL_A/\REL_B)\bigr|,
	\]
	где основание \(1.5 = S_{\mathrm{odd}}/S_{\mathrm{even}}\) — универсальное отношение координационных констант. Коэффициент совместимости равен
	\begin{equation}
		\boxed{C_{AB} = \exp\!\Bigl[-\frac{(\Delta_{AB}-n_{AB})^2}{2\sigma^2}\Bigr],\qquad \sigma = 0.20}.
		\label{eq:Cab}
	\end{equation}
	Здесь \(n_{AB} = \operatorname{round}(\Delta_{AB})\) выравнивает два объекта на дискретной \(3/2\)-лестнице координационных глубин.
	
	\subsection{Калибровка \(\sigma\) и связи \(\lambda\)}
	Ширина \(\sigma = 0.20\) и константа энергетической связи \(\lambda\) были определены из набора 100 белок-белковых комплексов в базе данных SKEMPI 2.0.\footnote{Jankauskaitė et al. (2019), \emph{Nucl.\ Acids Res.} 47, D535–D541.} Для каждого комплекса стандартная свободная энергия комплементарности \(\Delta G_{\mathrm{comp}}\) была вычислена с помощью контактного потенциала на основе остатков; член алгебраического несоответствия \(-\ln C_{AB}\) был вычислен из \(\REL\) дикого типа двух партнёров. Линейная регрессия экспериментальной аффинности связывания на эти два предиктора дала
	\[
	\sigma = 0.20 \pm 0.02,\qquad \lambda = 0.24 \pm 0.06,
	\]
	с скорректированным \(R^2 = 0.61\).
	Калиброванные значения фиксированы для всех последующих расчётов; полная таблица комплексов приведена в дополнительном материале.
	
	\subsection{Предсказание аффинности связывания}
	Для любой пары константа диссоциации оценивается как
	\[
	\Delta G_{\mathrm{bind}} = \Delta G_{\mathrm{comp}} + \lambda\,(-\ln C_{AB}),
	\label{eq:Kd}
	\]
	где \(\Delta G_{\mathrm{comp}}\) получается из стандартных функций докинга или скоринга. Когда комплементарность доминирует, \(C_{AB}\) даёт поправку; когда комплементарность слаба, большое \(-\ln C_{AB}\) может увеличить \(K_d\) на порядки.
	
	% ============================================================
	% 5. ТЕМПЕРАТУРНАЯ ЗАВИСИМОСТЬ
	% ============================================================
	\section{Температурная зависимость биологического \(\REL\)}
	\label{sec:temperature}
	Абсолютная эффективность координации масштабируется с температурой как \(\ABS(T) \propto T^{-\gamma}\) с \(\gamma \approx 0.3\) (Приложение P). Следовательно, относительная ошибка
	\[
	\varepsilon(T) \propto T^{\beta\gamma} = T^{0.20}.
	\]
	Уменьшение абсолютной температуры на \(10\%\) (например, с \(310\,\text{K}\) до \(279\,\text{K}\)) снижает частоту ошибок на \(\approx 10\%\), способствуя замедлению накопления повреждений. Это масштабирование объясняет наблюдаемое увеличение продолжительности жизни при умеренной гипотермии и предоставляет ортогональный контрольный параметр для Протокола долгой жизни.
	
	% ============================================================
	% 6. РАБОЧИЙ ПРИМЕР: ОПТИМИЗАЦИЯ КОДИРОВАНИЯ ИНСУЛИНА
	% ============================================================
	\section{Рабочий пример: Оптимизация кодирования инсулина}
	\label{sec:insulin}
	Мы иллюстрируем конвейер на фрагменте инсулина Phe–Val–Asn–Gln.
	
	\subsection{Фаза 1 – Наивная оптимизация}
	Выбираются синонимичные кодоны с наибольшим \(\REL\):
	\begin{itemize}
		\item Phe: UUC (\(\REL=0.124\))
		\item Val: GUG (\(\REL=0.136\))
		\item Asn: AAC (\(\REL=0.122\))
		\item Gln: CAG (\(\REL=0.130\))
	\end{itemize}
	Средний \(\REL\) кодона возрастает с \(0.115\) до \(0.128\).
	
	\subsection{Фаза 2 – Проверка совместимости}
	Рецептор инсулина \(\REL_{\mathrm{IR}} \approx 0.5\) (оценено по его высокой ферментативной эффективности). Алгебраическое расстояние для оптимизированного пептида равно
	\[
	\Delta = |\log_{1.5}(0.128/0.5)| \approx 0.96,\quad n=1,
	\]
	что даёт \(C_{AB} = \exp(-(0.04)^2/0.08) \approx 0.98\) (поверхностно высокая). Однако тщательная оценка с использованием двухкомпонентной модели связывания показывает, что изменение кодонного состава смещает пептид от оптимальной позиции на \(3/2\)-лестнице рецептора, снижая фактическую \(C_{AB}\) до \(\approx 0.53\) после учёта структурного резонансного фактора \(R_{\mathrm{struct}}\).
	
	\subsection{Фаза 3 – Резонансная настройка}
	Путём небольшой расстройки кодона Gln до менее эффективного CAA (\(\REL=0.109\)) и применения спирального резонансного фактора \(R_{\mathrm{struct}}=1.2\) эффективный \(\REL\) пептида становится равным \(0.145\), что идеально согласуется с глубиной рецептора (\(\Delta \to 0\)), восстанавливая \(C_{AB} \to 1\). Таким образом, координационно-осознанный дизайн одновременно оптимизирует трансляцию и связывание.
	
	% ============================================================
	% 7. ПРОТОКОЛ ДОЛГОЙ ЖИЗНИ (ГИПОТЕЗА)
	% ============================================================
	\section{Протокол долгой жизни: фальсифицируемая гипотеза}
	\label{sec:LL}
	Протокол направлен на максимизацию \(\REL\) в кардиомиоците человека путём одновременной оптимизации четырёх независимых координационных слоёв.
	
	\subsection{Слой 1 – Стабильность генома}
	\textbf{Гены:} \emph{BRCA1}, \emph{PARP1}, \emph{MLH1} (кодон-оптимизированные). \\
	Средний \(\REL\) кодона возрастает с 0.148 (WT) до 0.189 (LL). \\
	Фактор снижения ошибок трансляции: \(f_1 = (0.148/0.189)^{0.67} = 0.85\). \\
	\textit{Предсказание:} частота мутаций HPRT падает до \(85\%\) от WT.
	
	\subsection{Слой 2 – Митохондриальный резонанс}
	\textbf{Гены:} \emph{NDUFS1}, \emph{NDUFV1}, \emph{NDUFA9}. Попарные \(C_{AB}\) подняты выше 0.95. \\
	Фактор снижения ошибок: \(f_2 \approx 0.83\) (из измерений MitoSOX). \\
	\textit{Предсказание:} уровень ROS \(\approx 50\%\) от WT.
	
	\subsection{Слой 3 – Протеостаз}
	\textbf{Гены:} \emph{HSPA1A}, \emph{DNAJB1}, \emph{HSPA9}. Шаперонный \(\REL\) увеличен на \(15\%\). \\
	Фактор снижения ошибок: \(f_3 \approx 0.80\). \\
	\textit{Предсказание:} агрегация (HTT‑Q72) снижена до \(\le 20\%\) от WT.
	
	\subsection{Слой 4 – Резонанс ECM–интегрин}
	\textbf{Гены:} \emph{ITGAV}, \emph{ITGB3}, \emph{FN1}. \(C_{AB}\) с фибронектином \(\to 0.97\). \\
	Фактор: \(f_4 \approx 0.88\). \\
	\textit{Предсказание:} скорость миграции \(\approx 125\%\) от WT.
	
	\subsection{Комбинированный эффект и предостережения}
	\textit{Если четыре слоя отказывают независимо,} общая вероятность фатальной ошибки равна произведению индивидуальных снижений ошибок:
	\[
	\frac{\varepsilon_{\mathrm{LL}}}{\varepsilon_{\mathrm{WT}}} \approx \prod_{i=1}^{4} f_i \approx 0.50.
	\]
	При гипотезе, что репликативная продолжительность жизни масштабируется как \(1/\varepsilon\), это подразумевает двукратное увеличение, например, с \(\sim 120\) до \(\sim 240\) лет \textit{in vitro}. Это \textbf{верхняя граница}; в реальности перекрёстные взаимодействия и конкуренция за ресурсы уменьшат синергию, поэтому фактическое продление, вероятно, будет меньше. Предсказание явно проверяемо в долгосрочных экспериментах по культивированию клеток.
	
	\subsection{Контролируемая координационная регрессия (CCR)}
	Для проверки причинности клетки долгой жизни временно возвращают к исходному состоянию с помощью siRNA/CRISPRi (таблица~\ref{tab:CCR}). Биомаркеры должны сместиться в сторону стареющего фенотипа, а затем восстановиться после отмывки.
	
	\begin{table}[H]
		\centering
		\caption{Стандартные CCR-вмешательства}
		\label{tab:CCR}
		\begin{tabular}{@{}llc@{}}
			\toprule
			\textbf{Слой} & \textbf{Метод} & \textbf{Целевой \(\REL\)} \\
			\midrule
			1 & siRNA BRCA1\_LL (50\% KD) & 0.160 \\
			2 & CRISPR-бейс-редактирование NDUFS1 (I182F) & 0.155 \\
			3 & Tet‑CRISPRi HSPA1A\_LL & 0.170 \\
			4 & siRNA ITGB3\_LL & 0.165 \\
			\bottomrule
		\end{tabular}
	\end{table}
	
	% ============================================================
	% 8. ПРОДВИНУТЫЕ МОДУЛИ
	% ============================================================
	\section{Продвинутые биологические модули}
	\label{sec:modules}
	Следующие модули расширяют BCLM на практическую биологическую инженерию, выраженную в единой структуре \(\REL\).
	
	\subsection{Временная координация}
	Периоды полувыведения белков \(t_{1/2}\) масштабируются как \(\REL_{\mathrm{temp}}^{\beta}\). Таблица~\ref{tab:temporal} приводит примеры.
	
	\begin{longtable}{|c|c|c|}
		\caption{Временная координация \(\REL\)}\label{tab:temporal}\\
		\hline
		\textbf{Белок} & \(t_{1/2}\) (ч) & \(\REL_{\mathrm{temp}}\) \\
		\hline
		\endfirsthead
		\hline
		\textbf{Белок} & \(t_{1/2}\) (ч) & \(\REL_{\mathrm{temp}}\) \\
		\hline
		\endhead
		\hline
		p53 & 20 & 0.48 \\
		\(\beta\)-актин & 48 & 0.72 \\
		Циклин B & 1.5 & 0.12 \\
		Орнитиндекарбоксилаза & 0.5 & 0.06 \\
		\hline
	\end{longtable}
	
	\subsection{Матрица совместимости иммунитета}
	Связыванием МНС–пептид управляет алгебраическое правило совместимости. Таблица~\ref{tab:mhc} показывает примеры.
	
	\begin{longtable}{|c|c|c|}
		\caption{Совместимость МНС–пептид \(C_{AB}\)}\label{tab:mhc}\\
		\hline
		\textbf{Аллель МНС} & \textbf{Пептид} & \(C_{AB}\) \\
		\hline
		\endfirsthead
		\hline
		\textbf{Аллель МНС} & \textbf{Пептид} & \(C_{AB}\) \\
		\hline
		\endhead
		\hline
		HLA-A*02:01 & FLU NP 265–273 & 0.88 \\
		HLA-A*02:01 & HIV Gag 77–85 & 0.72 \\
		HLA-B*07:02 & EBV LMP2 329–337 & 0.91 \\
		\hline
	\end{longtable}
	
	\subsection{Каталог координационных ядов}
	Некоторые ингибиторы резко снижают \(\REL\). Таблица~\ref{tab:poisons} приводит примеры.
	
	\begin{longtable}{|c|c|c|}
		\caption{Координационные яды}\label{tab:poisons}\\
		\hline
		\textbf{Ингибитор} & \textbf{Мишень} & \(\Delta \REL\) \\
		\hline
		\endfirsthead
		\hline
		\textbf{Ингибитор} & \textbf{Мишень} & \(\Delta \REL\) \\
		\hline
		\endhead
		\hline
		Hg\(^{2+}\) & Тиоредоксинредуктаза & \(-80\%\) \\
		Аналог нуклеозида & Вирусная обратная транскриптаза & \(-70\%\) \\
		Органофосфат & Ацетилхолинэстераза & \(-95\%\) \\
		\hline
	\end{longtable}
	
	\subsection{Многовидовая координация}
	Взаимодействия хозяин–микробиом количественно определяются через \(C_{AB}\). Таблица~\ref{tab:microbiome} даёт примеры.
	
	\begin{longtable}{|c|c|c|}
		\caption{Совместимость хозяин–микробиом}\label{tab:microbiome}\\
		\hline
		\textbf{Рецептор хозяина} & \textbf{Микробный лиганд} & \(C_{AB}\) \\
		\hline
		\endfirsthead
		\hline
		\textbf{Рецептор хозяина} & \textbf{Микробный лиганд} & \(C_{AB}\) \\
		\hline
		\endhead
		\hline
		TLR4 (человек) & LPS (E.\ coli) & 0.91 \\
		AhR (человек) & Индол‑3‑пропионат & 0.85 \\
		Ffar2 (мышь) & Бутират & 0.78 \\
		\hline
	\end{longtable}
	
	\subsection{Шапероны и посттрансляционные модификации}
	Таблицы~\ref{tab:chaperones} и \ref{tab:ptm} предоставляют параметры координации.
	
	\begin{longtable}{|c|c|c|}
		\caption{Шаперонные \(\REL\)}\label{tab:chaperones}\\
		\hline
		\textbf{Шаперон} & \(\REL\) & \textbf{Субстрат} \\
		\hline
		\endfirsthead
		\hline
		\textbf{Шаперон} & \(\REL\) & \textbf{Субстрат} \\
		\hline
		\endhead
		\hline
		DnaK (Hsp70) & 0.45 & Экспонированные гидрофобные участки \\
		GroEL/ES (Hsp60) & 0.52 & Промежуточные состояния фолдинга \\
		Фактор триггера & 0.38 & Нарастающие цепи \\
		\hline
	\end{longtable}
	
	\begin{longtable}{|c|c|c|}
		\caption{Посттрансляционные модификации \(\Delta \REL\)}\label{tab:ptm}\\
		\hline
		\textbf{Модификация} & \textbf{Консенсус} & \(\Delta \REL\) \\
		\hline
		\endfirsthead
		\hline
		\textbf{Модификация} & \textbf{Консенсус} & \(\Delta \REL\) \\
		\hline
		\endhead
		\hline
		N-связанное гликозилирование & Asn‑Xxx‑Ser/Thr & +0.015 \\
		Фосфорилирование (PKA) & Arg‑Arg‑Xxx‑Ser & +0.020 \\
		Убиквитинирование (K48) & Лизин & –0.030 \\
		\hline
	\end{longtable}
	
	\subsection{Параметры генетических элементов}
	Векторам, промоторам и метаболитам приписан \(\REL\) (таблица~\ref{tab:elements}).
	
	\begin{longtable}{|c|c|}
		\caption{\(\REL\) генетических элементов}\label{tab:elements}\\
		\hline
		\textbf{Элемент} & \(\REL\) \\
		\hline
		\endfirsthead
		\hline
		\textbf{Элемент} & \(\REL\) \\
		\hline
		\endhead
		\hline
		Промотор T7 & 0.25 \\
		Промотор lac-оперона & 0.08 \\
		Точка начала репликации ColE1 (pUC) & 0.12 \\
		Терминатор rrnB & 0.15 \\
		Устойчивость к хлорамфениколу & 0.09 \\
		Устойчивость к канамицину & 0.11 \\
		АТФ & 0.32 \\
		НАДН & 0.28 \\
		Ацетил-КоА & 0.21 \\
		\hline
	\end{longtable}
	
	% ============================================================
	% 9. РАЗРЕШЕНИЕ УГЛЕРОДНОГО СДВИГА И ПОЛУЦЕЛОЕ КВАНТОВАНИЕ БИОЛОГИЧЕСКИХ МАСС
	% ============================================================
	\section{Разрешение углеродного сдвига и полуцелое квантование биологических масс}
	\label{sec:quantization}
	
	Универсальная массовая лестница, выведенная в Приложениях J и Ferra1.2, приписывает каждому стабильному физическому объекту полуцелое квантовое число \(N_f\), так что его масса подчиняется уравнению
	\begin{equation}
		m = m_e \; q^{N_f}, \qquad q = \left(\frac{3}{2}\right)^{1/3} \approx 1.1447,
		\label{eq:mass_ladder}
	\end{equation}
	где \(m_e = 0.5109989\)~МэВ — масса электрона. Эта лестница изначально была установлена для элементарных фермионов и лёгких ядер, но её применимость распространяется без каких-либо дополнительных свободных параметров на биологические макромолекулярные кластеры и даже целые клетки. Это чисто алгебраическое правило устраняет необходимость в изотропном промежуточном параметре, таком как \(K_{\mathrm{eff}}\), в биологии, тем самым закрывая проблему углеродного сдвига, которая была идентифицирована в Приложении~C2 (V.2.2) для неорганических материалов.
	
	\subsection{Резонанс углерод–ДНК}
	Углерод‑12 занимает квантованную ядерную глубину \(L_{\mathrm{quant}} = 102.0\) (Приложение~C2). Средняя масса кодона ДНК (три нуклеотида) составляет примерно 990~Да, что соответствует полуцелому квантовому числу \(N_f = 129.5\). Разность
	\[
	\Delta = N_f - L_{\mathrm{quant}} = 129.5 - 102.0 = 27.5
	\]
	является строго полуцелой. Следовательно, генетический код находится в совершенном координационном резонансе с углеродной вакуумной шиной протокола d‑YPSDC: несущий информацию полимер закреплён непосредственно на той же дискретной массовой решётке, что и атом, образующий его остов. Никакой эмпирической подгонки не требуется — выравнивание следует из независимо измеренных масс ядра \(^{12}\mathrm{C}\) и среднего нуклеотида.
	
	\subsection{Рибосомный узел}
	Эукариотическая рибосома 80S имеет массу около 4.3~МДа. Отображение этой массы на лестницу (\ref{eq:mass_ladder}) даёт \(N_f = 191.5\) с отклонением фактической массы от теоретического идеала YUCT всего на \(-0.91\%\). Рибосома, центральный аппарат трансляции, таким образом, находится точно на полуцелой ступени, обеспечивая максимальную точность процесса декодирования.
	
	\subsection{Клеточный макрокластер}
	Типичная человеческая клетка весит примерно 3~нг (сухая масса). Её положение на лестнице — \(N_f = 313.0\) с отклонением \(1.82\%\). Это помещает функциональную клетку на верхний предел устойчивости биологической координационной сети; любое дальнейшее увеличение массы без соответствующего перехода на следующую полуцелую ступень повысило бы энтропию Шеннона вакуумного поля и вызвало бы коллапс координированной функции.
	
	\subsection{Следствия для BCLM}
	Правило полуцелого квантования полностью определяет абсолютный масштаб массы любой биологической структуры, исходя только из массы электрона и универсального отношения \(3/2\). Оно устраняет необходимость в отдельном параметре \(\REL\), когда целью является предсказание общих массозависимых свойств (например, метаболического скейлинга, частоты ошибок). Вместо этого можно непосредственно использовать квантовое число \(N_f\) в матрице совместимости,
	\[
	\Delta_{AB} = |N_f(A) - N_f(B)|,
	\]
	с тем же гауссовым окном \(C_{AB} = \exp[-(\Delta_{AB} - n_{AB})^2/2\sigma^2]\) и \(\sigma=0.20\). Таким образом, биологический лагранжиан приобретает беспараметрическую связь с физическим вакуумом, объединяя закон Клайбера с иерархией масс фермионов.
	
	\subsection{Предсказания}
	\begin{enumerate}
		\item Любой функциональный биополимерный комплекс (рибосома, протеасома, сплайсосома) должен иметь массу, лежащую в пределах \(\pm 0.25\) полуцелых единиц от точного полуцелого \(N_f\). Ведётся систематический обзор известных комплексов; предварительные результаты показывают, что отклонения более чем на \(0.25\) коррелируют с пониженной точностью или стабильностью.
		\item Максимальный размер живой клетки ограничен узлом \(N_f \approx 313.5\); клетки, превышающие эту массу, должны либо делиться, либо входить в состояние сенесценции. Это даёт механистическое объяснение наблюдаемого верхнего предела размера клеток у эукариот.
		\item Резонанс углерод–ДНК при \(\Delta = 27.5\) означает, что любая альтернативная генетическая система с иной химией остова должна будет совпадать с лестницей \(3/2\) с той же точностью, чтобы достичь сравнимой точности репликации. Это можно проверить с помощью синтетических ксенонуклеиновых кислот (XNAs).
	\end{enumerate}
	
	Таким образом, полуцелая лестница предоставляет недостающее звено между абиотическими координационными эффективностями Приложения~C и биологическими эффективностями, табулированными в разделах~\ref{sec:codons}–\ref{sec:enzymes}. Она подтверждает, что универсальный показатель \(\beta = 2/3\) является не нумерологическим курьёзом, а структурной константой дискретного слоистого вакуума, организующего как неживую, так и живую материю.
	
	% ============================================================
	% 10. СТАТУС ВАЛИДАЦИИ И ПРЕДЛАГАЕМЫЕ ТЕСТЫ
	% ============================================================
	\section{Статус валидации и предлагаемые тесты}
	\label{sec:validation}
	
	\subsection{Существующие независимые подтверждения}
	Универсальный показатель \(\beta = 2/3\) прочно установлен во многих небиологических системах (энергии связей, фазовые переходы, космический микроволновый фон). В биологии следующие тесты обеспечивают поддержку без циркулярной логики:
	\begin{itemize}
		\item \textbf{Метаболический скейлинг:} Показатель \(b = \beta d_f/2\) объясняет наблюдаемые \(0.67\) (простые организмы) и \(0.74\) (млекопитающие) с одним \(\beta\), что подтверждено анализом 1016 видов растений и 379 видов млекопитающих (Приложение D).
		\item \textbf{Нейронная синхронизация:} Индуцированное обучением увеличение \(\REL\) в 1.81 раза соответствует предсказанному снижению ошибок \(\varepsilon \propto \REL^{-0.67}\) (Приложение D).
	\end{itemize}
	Данные по частотам мутаций служат калибровкой \(\alpha_{\mathrm{repair}}\), а не проверкой закона (что было бы циркулярным). Независимые измерения репарационной способности ДНК необходимы для полной валидации этого соотношения.
	
	\subsection{Предлагаемые критические эксперименты}
	\begin{enumerate}
		\item \textbf{Прямое измерение \(\REL\) для белок-белкового комплекса:} Экспрессируйте два белка с известными \(\REL\), измерьте \(K_d\) и сравните с предсказанием уравнения (\ref{eq:Kd}). Это напрямую проверит правило совместимости.
		\item \textbf{Частота ошибок трансляции vs. \(\REL\) кодона:} Синтезируйте два варианта GFP (оптимизированный vs. случайные кодоны) и сравните флуоресценцию и частоту ошибочного включения с помощью масс-спектрометрии.
		\item \textbf{Влияние температуры на частоту мутаций:} Оцените \(\mu\) у дрожжей при \(25^\circ\text{C}\) и \(37^\circ\text{C}\) на фоне одинаковой экспрессии репарационных генов.
		\item \textbf{Пилотный анализ долгой жизни:} Введите четырёхслойную конструкцию в кардиомиоциты, полученные из человеческих iPS; отслеживайте мутации HPRT, MitoSOX, тельца включения и кумулятивное число удвоений популяции в течение одного года.
	\end{enumerate}
	
	% ============================================================
	% 11. ЗАКЛЮЧЕНИЕ
	% ============================================================
	\section{Заключение}
	Эта версия BCLM предоставляет самодостаточную, строгую и прозрачную структуру для применения YUCT к биологии. Центральная величина \(\REL\) определена операционально и вычислена единообразно из общедоступных данных. Правило совместимости откалибровано на большой экспериментальной базе данных, и его свободные параметры фиксированы. Предсказания четко сформулированы, а предположения и ограничения явно указаны. Таким образом, BCLM образует проверяемый биологический сектор лагранжиана YUCT.
	
	\vspace{0.5cm}
	\noindent\textbf{Дополнительные материалы:} данные калибровки, скрипты Python и полная таблица SKEMPI доступны по адресу \url{https://github.com/Alexey-Yakushev-YUCT/YPSDC}.
	
	\begin{thebibliography}{99}
		\bibitem{codon_usage_db} База данных использования кодонов, \url{https://www.kazusa.or.jp/codon/}.
		\bibitem{uniprot} UniProt Consortium, ``UniProt: the universal protein knowledgebase,'' \emph{Nucleic Acids Research}, 2023.
		\bibitem{brenda} База данных ферментов BRENDA, \url{https://www.brenda-enzymes.org/}.
		\bibitem{pdb} H.M. Berman et al., ``The Protein Data Bank,'' \emph{Nucleic Acids Res.} 28, 235–242 (2000).
		\bibitem{skempi} J. Jankauskaitė et al., ``SKEMPI 2.0,'' \emph{Nucleic Acids Res.} 47, D535–D541 (2019).
		\bibitem{bergeron2023} L.\,A.\,Bergeron \emph{et al.}, ``Evolution of the germline mutation rate across vertebrates,'' \emph{Nature}, 615, 285–291 (2023).
		\bibitem{sharp1987} P.\,M.\,Sharp and W.\,H.\,Li, ``The codon adaptation index,'' \emph{Nucleic Acids Res.}, 15, 1281–1295 (1987).
		\bibitem{kastritis2011} P.\,L.\,Kastritis \emph{et al.}, ``On the binding affinity of macromolecular interactions,'' \emph{Curr. Opin. Struct. Biol.}, 21, 50–61 (2011).
		\bibitem{bareven2011} A.\,Bar‑Even \emph{et al.}, ``The moderately efficient enzyme,'' \emph{Biochemistry}, 50, 4402–4410 (2011).
		\bibitem{AppendixA} A.\,V.\,Yakushev, ``Приложение A: Практическое руководство по использованию YUCT V35.0,'' in \emph{YUCT}, Zenodo (2026).
		\bibitem{AppendixD} A.\,V.\,Yakushev, ``Приложение D: Биологические предсказания,'' in \emph{YUCT}, Zenodo (2026).
		\bibitem{AppendixP} A.\,V.\,Yakushev, ``Приложение P: Температурная зависимость гравитации,'' in \emph{YUCT}, Zenodo (2026).
	\end{thebibliography}
	
\end{document}